% !TeX program = xelatex
% 编译：xelatex -interaction=nonstopmode -halt-on-error main.tex（运行两遍）
% 封面信息在下方集中修改。缺少图片时自动显示占位框；添加同名文件后自动显示图片。
\documentclass[UTF8,a4paper,zihao=-4,fontset=fandol]{ctexart}
\usepackage[top=2.6cm,bottom=2.5cm,left=2.7cm,right=2.7cm,headheight=16pt]{geometry}
\usepackage{amsmath,amssymb,bm}
\usepackage{graphicx,booktabs,tabularx,array}
\usepackage{caption,float,placeins,fancyhdr,setspace,enumitem}
\usepackage[hidelinks]{hyperref}
\usepackage{xurl}
\usepackage{xparse}
\usepackage{needspace}
% Generated by report/generate_ppo_figures.py from measured scores.
\newcommand{\RandomMean}{0.88}
\newcommand{\RandomStd}{1.03}
\newcommand{\RandomMedian}{1.00}
\newcommand{\RandomPxxv}{0.00}
\newcommand{\RandomMax}{4.00}
\newcommand{\RandomLow}{50}
\newcommand{\RandomTimeout}{0}
\newcommand{\DQNMean}{273.36}
\newcommand{\DQNStd}{84.28}
\newcommand{\DQNMedian}{266.00}
\newcommand{\DQNPxxv}{235.50}
\newcommand{\DQNMax}{497.00}
\newcommand{\DQNLow}{0}
\newcommand{\DQNTimeout}{0}
\newcommand{\SourceMean}{129.98}
\newcommand{\SourceStd}{119.51}
\newcommand{\SourceMedian}{84.00}
\newcommand{\SourcePxxv}{41.00}
\newcommand{\SourceMax}{411.00}
\newcommand{\SourceLow}{18}
\newcommand{\SourceTimeout}{0}
\newcommand{\PPOControlMean}{196.54}
\newcommand{\PPOControlStd}{138.91}
\newcommand{\PPOControlMedian}{166.50}
\newcommand{\PPOControlPxxv}{80.00}
\newcommand{\PPOControlMax}{436.00}
\newcommand{\PPOControlLow}{9}
\newcommand{\PPOControlTimeout}{5}
\newcommand{\PPOGaeMean}{203.42}
\newcommand{\PPOGaeStd}{126.34}
\newcommand{\PPOGaeMedian}{191.50}
\newcommand{\PPOGaePxxv}{69.75}
\newcommand{\PPOGaeMax}{410.00}
\newcommand{\PPOGaeLow}{6}
\newcommand{\PPOGaeTimeout}{3}
\newcommand{\DQNPPODiff}{76.82}
\newcommand{\DQNPPOCI}{[28.02,125.24]}
\newcommand{\DQNPPOWins}{31}
\newcommand{\DQNPPOTies}{0}
\newcommand{\DQNPPOLosses}{19}
\newcommand{\GaeControlDiff}{6.88}
\newcommand{\GaeControlCI}{[-41.08,54.68]}
\newcommand{\GaeControlWins}{25}
\newcommand{\GaeControlTies}{0}
\newcommand{\GaeControlLosses}{25}

% Generated by report/generate_ppo_analysis.py from measured data.
\newcommand{\TricksAblationDiff}{-9.67}
\newcommand{\TricksAblationLow}{-42.63}
\newcommand{\TricksAblationHigh}{20.87}
\newcommand{\ShapingAblationDiff}{-31.20}
\newcommand{\ShapingAblationLow}{-62.56}
\newcommand{\ShapingAblationHigh}{-1.82}
\newcommand{\BalancedAblationDiff}{18.92}
\newcommand{\BalancedAblationLow}{-13.52}
\newcommand{\BalancedAblationHigh}{52.16}
\newcommand{\DualRateAblationDiff}{-5.88}
\newcommand{\DualRateAblationLow}{-49.06}
\newcommand{\DualRateAblationHigh}{36.60}
\newcommand{\GaeAblationDiff}{6.88}
\newcommand{\GaeAblationLow}{-41.08}
\newcommand{\GaeAblationHigh}{54.68}
\newcommand{\KLAblationDiff}{7.81}
\newcommand{\KLAblationLow}{-47.29}
\newcommand{\KLAblationHigh}{64.17}
\newcommand{\KLRepeatOneDiff}{53.60}
\newcommand{\KLRepeatOneAdaptiveMean}{224.26}
\newcommand{\KLRepeatOneControlMean}{170.66}
\newcommand{\KLRepeatTwoDiff}{-33.50}
\newcommand{\KLRepeatTwoAdaptiveMean}{181.54}
\newcommand{\KLRepeatTwoControlMean}{215.04}
\newcommand{\KLRepeatThreeDiff}{3.34}
\newcommand{\KLRepeatThreeAdaptiveMean}{213.78}
\newcommand{\KLRepeatThreeControlMean}{210.44}
\newcommand{\KLAdaptiveMean}{206.53}
\newcommand{\KLControlMean}{198.71}

\newcommand{\CourseName}{强化学习}
\newcommand{\ReportTitle}{基于 DQN 与 PPO 的 Breakout 智能体设计与比较}
\newcommand{\MajorClass}{硕2605班}
\newcommand{\StudentID}{M202674292}
\newcommand{\StudentName}{赵恒升}
\newcommand{\Instructor}{何琨}
\newcommand{\ReportDate}{2026 年 10 月 22 日}
\newcommand{\CollegeName}{计算机科学与技术学院}
\newcolumntype{Y}{>{\raggedright\arraybackslash}X}
\newcolumntype{P}[1]{>{\raggedright\arraybackslash}p{#1}}
\setstretch{1.25}
\setlength{\parindent}{2em}
\setlength{\parskip}{0pt}
\setlength{\emergencystretch}{2em}
\raggedbottom
% 允许图表与正文共同填充页面，避免浮动比例限制产生半空页。
\setcounter{topnumber}{3}
\setcounter{bottomnumber}{2}
\setcounter{totalnumber}{4}
\renewcommand{\topfraction}{.92}
\renewcommand{\bottomfraction}{.85}
\renewcommand{\textfraction}{.06}
\renewcommand{\floatpagefraction}{.8}
\setlength{\textfloatsep}{12pt plus 2pt minus 2pt}
\setlength{\floatsep}{10pt plus 2pt minus 2pt}
\setlength{\intextsep}{10pt plus 2pt minus 2pt}
\widowpenalty=10000
\clubpenalty=10000
\renewcommand{\arraystretch}{1.15}
\setlist[enumerate]{leftmargin=2.3em,itemsep=2pt,topsep=4pt}
\captionsetup{font=small,labelsep=quad,skip=7pt}
\ctexset{
  section={format=\Large\heiti,beforeskip=1.2em,afterskip=.7em},
  subsection={format=\large\heiti,beforeskip=.9em,afterskip=.4em},
  subsubsection={format=\normalsize\heiti,beforeskip=.7em,afterskip=.3em},
  contentsname={目\quad 录}
}
\pagestyle{fancy}
\fancyhf{}
\fancyhead[C]{\small 华中科技大学课程实验报告}
\fancyfoot[C]{\small\thepage}
\renewcommand{\headrulewidth}{0.4pt}
\renewcommand{\footrulewidth}{0pt}
\fancypagestyle{plain}{%
  \fancyhf{}
  \fancyhead[C]{\small 华中科技大学课程实验报告}
  \fancyfoot[C]{\small\thepage}
  \renewcommand{\headrulewidth}{0.4pt}
  \renewcommand{\footrulewidth}{0pt}
}
\setcounter{tocdepth}{2}
\numberwithin{equation}{section}
\hypersetup{pdftitle={基于DQN与PPO的Breakout智能体课程实验报告},pdfsubject={强化学习课程项目}}

% \reportfigure[图片尺寸选项][浮动位置]{文件路径}{占位框标题}{占位提示}{图注}{标签}
% 替换时只需把图片放到相应路径，无须修改正文和交叉引用。
\NewDocumentCommand{\reportfigure}{O{width=.98\linewidth,height=14cm,keepaspectratio} O{!htbp} m m m m m}{%
  \begin{figure}[#2]
    \centering
    \IfFileExists{#3}{%
      \includegraphics[#1]{#3}
    }{%
      \fbox{\begin{minipage}[c][4.2cm][c]{.88\linewidth}
        \centering
        {\heiti #4}\par\vspace{.65em}
        {\small #5}\par\vspace{.8em}
        {\footnotesize 待插入文件：\nolinkurl{#3}}
      \end{minipage}}
    }
    \caption{#6}\label{#7}
  \end{figure}%
}
\newcommand{\coverline}[2]{%
  \makebox[3.0cm][l]{#1}\hspace{.35cm}
  \underline{\makebox[6.7cm][c]{#2}}\par\vspace{.55cm}
}

\begin{document}
\begin{titlepage}
  \thispagestyle{empty}
  \centering
  \vspace*{1.7cm}
  {\kaishu\fontsize{38}{48}\selectfont 华中科技大学\par}
  \vspace{.65cm}
  {\kaishu\fontsize{34}{44}\selectfont 课 程 实 验 报 告\par}
  \vspace{1.35cm}
  {\heiti\zihao{2} 课程名称：\ \underline{\makebox[6.1cm]{\CourseName}}\par}
  \vspace{.85cm}
  {\heiti\zihao{3}\ReportTitle\par}
  \vspace{2.45cm}
  \begin{minipage}{10.3cm}
    \centering\songti\bfseries\zihao{3}
    \coverline{专业班级}{\MajorClass}
    \coverline{学\qquad 号}{\StudentID}
    \coverline{姓\qquad 名}{\StudentName}
    \coverline{指导教师}{\Instructor}
    \coverline{报告日期}{\ReportDate}
  \end{minipage}
  \vfill
  {\songti\bfseries\zihao{3}\CollegeName\par}
  \vspace*{.65cm}
\end{titlepage}

\pagenumbering{Roman}
\tableofcontents
\clearpage
\pagenumbering{arabic}

\begin{center}
  {\Large\heiti \ReportTitle}
\end{center}
% 大章节独立起页；章节内部由正文与浮动图表自然分页。
\section{项目背景与目标}
\label{sec:introduction}
\subsection{研究背景与任务目标}
Breakout 的动作空间很小，但控制过程同时包含视觉感知、延迟奖励与连续决策。挡板移动通常不会立即得分；一次接球的收益可能在小球多次飞行和碰撞后才体现。智能体需要从连续图像中判断运动趋势，并将后续击砖收益与先前动作联系起来。

本项目在统一的自制固定球速 Breakout 环境中，实现 DQN 和 PPO 两条图像控制链路，研究动作价值学习与直接策略优化的原理、实现和训练表现。任务要求智能体仅根据图像选择四个离散动作，在完整五生命游戏中提高原生得分。环境规则由项目代码定义，结果对应这一任务，不直接对照公开 Atari 基准。

\subsection{研究内容与报告安排}
项目工作包括环境与输入接口建设、两种算法的训练与恢复、开发集模型选择、固定模型测试，以及网页实时展示。DQN 重点研究低学习率续训、训练退化与恢复实验；PPO 重点研究策略与价值网络、更新稳定性及 GAE 单因素对照。实验同时保留有效结果与负结果，以日志、检查点和逐局分数支持分析。

第\ref{sec:shared}节给出共用环境、系统接口和评价口径；第\ref{sec:dqn}节和第\ref{sec:ppo}节分别完整介绍 DQN、PPO 的原理、实现、训练及开发分析；第\ref{sec:comparison}节集中报告冻结模型的统一测试与方法比较；第\ref{sec:conclusion}节总结结论和局限。开发数据用于选模，最终测试用于评价已冻结的候选，二者分别报告。

\FloatBarrier

\clearpage
\section{游戏环境与公共实验基础}
\label{sec:shared}
\subsection{游戏规则与物理过程}
项目在 Gymnasium 接口下实现固定球速 Breakout 环境。画面分辨率为 $210\times160$，包含小球、挡板、砖块、边界以及分数和剩余生命显示。每局有 5 条生命，每面砖墙有 $6\times18=108$ 块砖，共设置两面砖墙。各行砖块从上至下分别计 7、7、4、4、1、1 分，因此一面墙为 432 分，整局最高分为 864 分。

物理系统以 60 Hz 更新，小球每帧移动 2 像素，即球速固定为 120 像素/秒。挡板宽度固定为 16 像素，单个物理帧向左或向右移动 3 像素。碰撞后小球改变运动方向但保持速度模长。将球速和挡板宽度固定，可以减少随游戏进程变化的额外难度因素，使模型在一致的控制条件下学习。

小球与挡板接触的位置决定反弹方向。挡板横向划分为五个区域，越靠近边缘，反弹方向越偏向相应侧；中心区域保留小角度水平分量，避免长时间完全垂直往返。该规则不随累计接球次数改变。砖块碰撞使用运动路径与扩展矩形的相交检测，在一个物理步内寻找最早接触点，处理反射并继续推进剩余时间，以避免仅检查离散终点造成漏碰撞。

发球位置和初始水平方向由随机种子控制。第一面砖墙清空后重建第二面墙并重新发球；生命耗尽或第二面砖墙清空时游戏真正结束。该环境是按任务需要重建的游戏，其几何和反弹规则由项目代码定义，实验分数对应这一固定规则环境。

\reportfigure[width=.72\linewidth,height=9cm,keepaspectratio]{images/game_live.png}{游戏实时展示画面}{网页运行截图，展示游戏区域和演示控制。}{DQN 智能体在共用网页展示接口中的实时运行画面，显示游戏区域、得分、剩余生命与演示控制。}{fig:game}

\subsection{状态、动作与奖励定义}
物理环境的完整状态可以表示为小球位置与速度、挡板位置、砖块存活标记、剩余生命、当前墙面以及发球状态等变量的集合。这些变量决定下一步物理变化。智能体实际获得的是渲染图像，模型不直接读取球坐标、速度或预先计算的落点。

动作空间为四个离散动作，如表\ref{tab:actions}所示。自动发球包装器在等待发球时将当前动作替换为发球动作，并在正常的物理步中执行，不额外插入隐藏帧。因此，学习任务主要集中在接球与击砖控制。

\begin{table}[!htbp]
 \centering
 \caption{环境动作空间与控制含义}\label{tab:actions}
 \begin{tabularx}{.86\linewidth}{l Y}
 \toprule
 动作名称 & 物理控制含义\\
 \midrule
 NOOP & 挡板保持当前位置\\
 FIRE & 等待发球时发射小球\\
 RIGHT & 挡板向右移动，受右边界约束\\
 LEFT & 挡板向左移动，受左边界约束\\
 \bottomrule
 \end{tabularx}
\end{table}

奖励直接采用当前动作执行期间击碎砖块获得的原生分数：
\begin{equation}
 r_t=\sum_{b\in\mathcal B_t}v(b),
\end{equation}
其中 $\mathcal B_t$ 为该决策步内被击碎的砖块集合，$v(b)$ 为砖块所在行对应的分值。DQN 和最终 PPO 对照均采用原始奖励，不设置接球奖励、存活奖励或失球惩罚，也不裁剪奖励。辅助奖励的单独消融见第\ref{sec:ppo}节。所有评估只累计原生得分。失球虽无直接负奖励，但会减少剩余生命和未来得分机会。

\subsection{图像预处理与时序信息}
单张画面可以反映小球的位置，却无法可靠区分运动方向。项目将连续四张处理后的画面堆叠，作为对运动状态的近似表示。共享环境工厂按以下顺序处理图像。

首先，将一个动作重复执行 4 个物理帧，并累加期间的奖励。随后，对实际产生的最后两帧逐像素取最大值，得到池化画面。若提前终止或截断，则使用真正发生的末尾画面，避免把旧帧当成终止观测。之后，使用灰度转换和缩放得到 $84\times84$ 图像，再按时间顺序堆叠 4 帧，并将通道维移动到最前端。

因此，单个状态表示为
\begin{equation}
 s_t=[o_{t-3},o_{t-2},o_{t-1},o_t]
 \in\{0,\ldots,255\}^{4\times84\times84}.
\end{equation}
单环境批输入的形状为 $(1,4,84,84)$，原始存储类型为 \texttt{uint8}。网络内部将像素转换为浮点数并除以 255。堆帧允许卷积网络利用位置变化判断运动趋势，但它仍是从图像得到的近似状态表示。

\reportfigure{images/observation_pipeline.png}{图像预处理示意}{原始 RGB 画面、最大池化结果与四帧灰度输入。}{共享环境工厂实际采集的一组连续观测。四张灰度图按时间顺序组成 $4\times84\times84$ 输入，最新一帧已逐像素核对实际网络观测。}{fig:preprocess}

\subsection{回合终止、时间截断与环境一致性}
项目以一整局游戏作为一个回合，损失一条生命不会立即终止回合。只要仍有剩余生命，环境就准备并自动发射下一球。这样可以让价值估计覆盖一局中的多次失球和后续得分机会。

所有模型使用统一的单局时间上限。达到上限属于时间截断，生命耗尽或通关属于真正终止，两者在学习目标中需要区分。真正终止时未来价值为零；时间截断则仍需要基于截断时的最终观测进行适当自举。环境包装器保留相应标记和终止观测，由训练库的回放与目标计算流程处理。

训练、独立评估、实时展示和视频录制共用同一套环境与预处理代码。模型保存环境版本和关键参数，加载时检查球速、挡板宽度、图像尺寸、堆帧数量和奖励规则等设置。该检查能够阻止外形相同但物理规则不同的历史模型被误用于当前任务。

\subsection{强化学习任务与长期回报}
设 $s_t$ 为时刻 $t$ 提供给智能体的状态表示，$a_t$ 为动作，$r_t$ 为执行该动作后获得的奖励，$T$ 为本局结束时刻。折扣累计回报定义为
\begin{equation}
 G_t=\sum_{k=0}^{T-t-1}\gamma^k r_{t+k},\qquad
 J(\pi)=\mathbb{E}_{\pi}[G_0],
\end{equation}
其中 $\pi$ 为策略，$\gamma\in[0,1)$ 为折扣因子。DQN 与初始 PPO 取 $\gamma=0.99$，后续 PPO 实验取 0.995，以调整对远期击砖收益的权重。训练目标是提高期望折扣回报；实验评估则报告整局未折扣的游戏总分，两者含义不同。

\subsection{共用系统结构与实时展示}
系统使用 Python、Gymnasium 和 Stable-Baselines3 实现 DQN 与 PPO。游戏、预处理、参数管理、模型加载、评估和实时推理接口共用。算法各自的采样、更新及续训机制分别在第\ref{sec:dqn}节和第\ref{sec:ppo}节说明。

环境层接收动作，推进物理过程并返回图像、奖励和结束标记；管理层定期评估并保存检查点；展示层用相同观测接口确定性推理。各层共享任务定义，训练和评估使用独立环境实例。

\begin{table}[!htbp]
 \centering\small
 \caption{共用程序模块及其职责}\label{tab:modules}
 \begin{tabularx}{\linewidth}{P{4.5cm}Y}
 \toprule
 模块 & 主要实现内容\\
 \midrule
 \texttt{fixed\_breakout.py} & 游戏物理、碰撞、计分、渲染、自动发球与动作重复\\
 \texttt{env.py} & 环境工厂、版本检查、灰度缩放、帧堆叠与通道转换\\
 \texttt{model\_io.py} & 识别 DQN/PPO 变体并检查游戏契约\\
 \texttt{callbacks.py} & 开发评估、最佳模型选择、早停与训练状态记录\\
 \texttt{evaluate.py} & 逐局种子评估、结果导出及随机策略比较入口\\
 \texttt{dashboard.py}、\texttt{game\_viewer.py} & 网页展示、模型切换、暂停及播放控制\\
 \bottomrule
 \end{tabularx}
\end{table}

展示采用确定性动作：DQN 选择最大 Q 值，PPO 选择最大动作概率。网页提供模型选择、随机策略、暂停、倍速和全屏，显示实际物理帧，1 倍以 60 帧/秒为目标。定量评价仍依赖多局统计。

\subsection{评价指标与数据使用原则}
每局游戏总分定义为 $R_i=\sum_{t=0}^{T_i-1}r_t$。对 $N$ 局游戏，平均分与总体形式的标准差分别为
\begin{equation}
 \overline R=\frac{1}{N}\sum_{i=1}^{N}R_i,\qquad
 \sigma_R=\sqrt{\frac{1}{N}\sum_{i=1}^{N}(R_i-\overline R)^2}.
\end{equation}
平均分衡量总体得分水平，中位数反映典型游戏表现，标准差和低分位数用于分析波动。仅报告最高分容易放大少数成功轨迹的影响；仅报告损失则不能直接反映控制效果。

模型选择使用开发数据，冻结候选后再汇总统一测试结果。相同开局种子上的逐局分数用于配对比较；区间按完整种子对进行 bootstrap 重采样。此类区间反映固定模型在回合样本上的差异，不覆盖独立重复训练的波动。不同算法的损失目标和尺度不同，损失只用于各自的训练诊断。

\begin{table}[!htbp]
 \centering\footnotesize
 \caption{报告中的数据用途与评估协议}\label{tab:data-protocols}
 \begin{tabularx}{\linewidth}{P{3.1cm}P{3.8cm}Y}
 \toprule
 数据用途 & 采样与隔离方式 & 报告位置与解释\\
 \midrule
 DQN 训练开发评估 & 固定随机性设置，连续生成游戏序列 & 第\ref{sec:dqn}节；用于各阶段选模\\
 DQN 开发复测 & 逐局重置，各候选共享开局集合 & 第\ref{sec:dqn}节；30 局候选诊断\\
 PPO 起点选择与开发 & 起点验证与训练开发使用不同集合 & 第\ref{sec:ppo}节；分别用于起点选择和训练选模\\
 历轮 PPO 实验评估 & 各轮使用各自记录的种子集合 & 第\ref{sec:ppo}节；不拼接为同测试集提升曲线\\
 冻结模型统一测试 & 使用留出的共同开局，逐局重置 & 第\ref{sec:comparison}节；50 局方法比较与最终 GAE 消融\\
 \bottomrule
 \end{tabularx}
\end{table}


\FloatBarrier

\clearpage
\section{DQN 方法与实验}
\label{sec:dqn}
本节按算法原理、模型实现、训练设置与开发分析的顺序介绍 DQN。除明确引用的统一测试外，本节分数均属于训练开发评估或开发复测。

\subsection{算法原理}
\subsubsection{动作价值与 Q-learning}
动作价值函数描述在状态 $s$ 下先执行动作 $a$，随后按策略 $\pi$ 行动时的预期回报：
\begin{equation}
 Q^\pi(s,a)=\mathbb{E}_{\pi}[G_t\mid s_t=s,a_t=a].
\end{equation}
在同一画面下，向左或向右移动可能带来不同的后续接球机会。价值函数将这种长期差异转化为可比较的数值。若最优动作价值函数 $Q^*$ 已知，就可以采用 $\arg\max_a Q^*(s,a)$ 选择动作。

最优动作价值满足 Bellman 最优关系，即当前动作的价值由即时奖励和下一状态的最优未来价值共同决定：
\begin{equation}
 Q^*(s,a)=\mathbb{E}\!\left[
 r+\gamma(1-d)\max_{a'}Q^*(s',a')\mid s,a
 \right],
\end{equation}
其中 $s'$ 是下一状态，$d$ 表示真正终止。表格型 Q-learning 利用一次环境转移更新对应表项：
\begin{equation}
 Q(s_t,a_t)\leftarrow Q(s_t,a_t)+\alpha
 \left[r_t+\gamma(1-d_t)\max_{a'}Q(s_{t+1},a')-Q(s_t,a_t)\right].
\end{equation}
这里 $\alpha$ 为学习率，方括号中的差值是时间差分误差。更新通过下一状态的估计值修正当前价值，因此不必等一局游戏全部结束后才开始学习。

图像状态的组合数量远大于表格能够存储的范围。DQN 用参数为 $\theta$ 的神经网络 $Q_\theta(s,a)$ 近似动作价值，并通过共享视觉特征将已见画面的经验迁移到相似画面。经典 DQN 将卷积表示与价值学习结合，建立了从像素输入到离散动作决策的学习框架。本项目沿用该框架，网络输出四个动作的 Q 值，而不是直接输出动作概率。

\subsubsection{经验回放与目标网络}
神经网络训练通常希望样本具有较低的相关性，而游戏中相邻画面非常相似。若严格按采样顺序立即更新，参数容易受近期轨迹支配。经验回放池 $\mathcal D$ 保存转移
\begin{equation}
 e_t=(s_t,a_t,r_t,s_{t+1},d_t),
\end{equation}
训练时从池中随机抽取小批量样本。这样既减弱相邻转移之间的相关性，也允许同一条有效经验被多次利用。低奖励和失球经验同样保留在池中，因为区分成功与失败对动作评价都具有意义。

DQN 同时维护在线网络 $Q_\theta$ 与目标网络 $Q_{\theta^-}$。在线网络参与梯度更新并用于决策；目标网络在一段时间内保持不变，用于构造学习目标，随后从在线网络复制参数。一步 TD 目标和误差写为
\begin{align}
 y_t &= r_t+\gamma(1-d_t)\max_{a'}Q_{\theta^-}(s_{t+1},a'),\\
 \delta_t &= y_t-Q_\theta(s_t,a_t).
\end{align}
目标网络将监督目标的变化与每次在线更新适度分离，缓解自举估计与参数更新相互追逐的问题。计算 $y_t$ 时停止梯度传播，梯度仅作用于在线网络。

本项目使用 Stable-Baselines3 的 DQN 实现，采用 Huber 损失：
\begin{equation}
 \ell(\delta)=
 \begin{cases}
 \frac{1}{2}\delta^2,& |\delta|\leq 1,\\
 |\delta|-\frac{1}{2},& |\delta|>1,
 \end{cases}
 \qquad
 \mathcal L(\theta)=\frac{1}{B}\sum_{i=1}^{B}\ell(\delta_i).
\end{equation}
其中 $B$ 为批大小。小误差区域使用平方形式，大误差区域使用线性形式，从而减弱个别较大 TD 误差对更新的影响。优化器采用 Adam。损失减小表示网络更接近当前样本上的 TD 目标，并不直接意味着下一次评估的游戏得分必然提高。

\subsubsection{探索机制与多步回报}
训练采用 $\varepsilon$-greedy 策略：
\begin{equation}
 a_t=
 \begin{cases}
 \text{从动作空间均匀采样},&\text{概率为 }\varepsilon_t,\\
 \arg\max_a Q_\theta(s_t,a),&\text{概率为 }1-\varepsilon_t.
 \end{cases}
\end{equation}
早期较大的 $\varepsilon_t$ 帮助收集不同状态与动作的经验，后期降低探索率以更多利用已学策略。开发评估采用贪心动作，因此训练阶段含随机探索的回合分数与评估分数应分别解释。

对于延迟奖励，一步目标需要通过多次自举将价值向前传播。项目还尝试了三步回报。中间没有真正终止时，其目标为
\begin{equation}
 y_t^{(3)}=r_t+\gamma r_{t+1}+\gamma^2r_{t+2}
 +\gamma^3\max_{a'}Q_{\theta^-}(s_{t+3},a').
\end{equation}
若提前终止，则仅累积终止前的奖励，并去除终止后的自举项。多步回报可以更直接地利用后续实际奖励，但也引入更长轨迹中的随机性和行为策略影响，其效果需要结合实验判断。本项目主要结果采用一步 DQN，三步设置作为后续恢复实验。

\subsection{模型实现与训练管理}
\subsubsection{卷积价值网络}
模型采用 \texttt{CnnPolicy}，使用 NatureCNN 风格的卷积特征提取器。输入的四帧图像沿通道维拼接，卷积同时提取空间结构与跨帧差异。三层卷积之后将特征展平，经过 512 维全连接层，再由动作价值输出层产生四个 Q 值。各层使用的尺寸列于表\ref{tab:network}。

\begin{table}[!htbp]
 \centering\small
 \caption{DQN 卷积价值网络结构，省略批维度}\label{tab:network}
 \begin{tabularx}{\linewidth}{l Y l}
 \toprule
 层 & 运算与参数 & 输出尺寸\\
 \midrule
 输入 & 四帧灰度图像；像素除以 255 & $4\times84\times84$\\
 卷积 1 & 32 个 $8\times8$ 卷积核，步长 4，ReLU & $32\times20\times20$\\
 卷积 2 & 64 个 $4\times4$ 卷积核，步长 2，ReLU & $64\times9\times9$\\
 卷积 3 & 64 个 $3\times3$ 卷积核，步长 1，ReLU & $64\times7\times7$\\
 展平 & 按通道和空间位置展开 & 3136\\
 全连接 & $3136\rightarrow512$，ReLU & 512\\
 输出层 & $512\rightarrow4$，线性输出 & 4\\
 \bottomrule
 \end{tabularx}
\end{table}

输出层不使用 softmax，因为 Q 值表示预期回报，可以是任意实数。动作选择依据输出值的相对大小进行。在线网络与目标网络具有相同结构，目标网络由在线网络周期性复制得到。当前实现使用普通 DQN 的最大目标值计算，没有加入 Double DQN、Dueling 结构或优先经验回放。

网络不显式预测小球落点，也不包含人工编写的追球控制规则。像素特征是否有利于接球，需要通过环境交互产生的奖励逐步学习。由此，模型实现将视觉表示与控制目标连接起来，而游戏的物理状态仅用于环境内部更新和调试。

\subsubsection{训练循环与模型更新}
DQN 训练流程如下，其中环境步均指经过动作重复包装后的决策步：
\begin{enumerate}
 \item 初始化在线网络，复制其参数得到目标网络，建立回放池并重置环境。
 \item 根据当前探索率选择动作，执行环境步，获得新观测、累计奖励和结束信息。
 \item 将转移存入回放池；回合结束后重置环境，继续收集完整游戏经验。
 \item 超过预热步数后，每收集 4 个决策步执行 1 次梯度更新，每次随机采样 32 条转移。
 \item 使用目标网络构造 TD 目标，以 Huber 损失更新在线网络；每 10000 个决策步同步目标网络。
 \item 定期运行独立开发评估，保存检查点和表现最好的模型，达到预算或早停条件后结束。
\end{enumerate}

这里动作重复 4 帧与每 4 个决策步训练一次是两个不同层次的设置。前者控制智能体多久重新作出一次决策，后者控制采样与参数更新的比例。在未提前结束的情况下，一次梯度更新通常对应最近约 16 个物理帧的采样。报告中的 50 万步等训练预算均采用决策步计数。

\subsubsection{续训与回放状态保存}
模型文件保存网络参数、优化器状态及相关配置，但不自动包含经验回放池。项目因此把回放池作为独立产物保存，并检查其环境定义、容量与回报跨度是否与待恢复模型兼容。压缩用于降低磁盘占用，训练时解压后的回放仍需要相应内存。

在缺少配套回放池的续训中，模型从已有权重开始建立新的回放池。项目实现了策略预热：前 20000 个决策步仍暂缓梯度更新，但动作由恢复后的 $\varepsilon$-greedy 策略产生。这样，预热数据与当前策略的访问状态更接近，而不是重新由完全随机动作填充。这种操作属于基于已有模型的新阶段训练，不等同于完整恢复原训练过程。

DQN 训练入口为 \texttt{train.py}，策略预热位于 \texttt{warm\_start\_dqn.py}，回放兼容性与压缩恢复由 \texttt{replay\_io.py} 实现。

\subsection{训练设置与日志诊断}
\subsubsection{基础训练设置}
DQN 基础实验在固定随机性设置下，从随机初始化开始训练 500000 个决策步。初始探索率为 1.0，在前 60\% 的训练预算内线性下降至 0.05，之后保持不变。学习率固定为 $10^{-4}$，回放容量为 50000，预热步数为 20000。主要参数与续训参数的对照见表\ref{tab:hyper}。

\begin{table}[!htbp]
 \centering\small
 \caption{从头训练与当前低学习率续训配置}\label{tab:hyper}
 \begin{tabularx}{\linewidth}{Y c c}
 \toprule
 参数 & 从头训练 & 低学习率续训\\
 \midrule
 阶段预算（决策步） & 500000 & 200000\\
 学习率 & $1\times10^{-4}$ & $3\times10^{-5}\rightarrow1\times10^{-5}$\\
 折扣因子 $\gamma$ & 0.99 & 0.99\\
 回放池容量 & 50000 & 50000\\
 新回放池预热步数 & 20000 & 20000\\
 批大小 & 32 & 32\\
 训练频率 / 每次梯度步数 & 4 / 1 & 4 / 1\\
 目标网络同步间隔 & 10000 & 10000\\
 探索率变化 & $1.00\rightarrow0.05$ & $0.05\rightarrow0.02$\\
 探索率衰减占阶段比例 & 60\% & 100\%\\
 回报跨度 & 1 步 & 1 步\\
 评估间隔（决策步） & 50000 & 25000\\
 每次评估局数 & 20 & 50\\
 \bottomrule
 \end{tabularx}
\end{table}

图像转移需要较大的内存。按当前回放规模估算，分别以 \texttt{uint8} 保存当前与下一堆帧观测时，仅观测约占 2.8 GB，其他字段和运行过程还会带来额外开销。因此，回放容量同时影响可保留的历史范围与内存需求。

训练过程记录回合奖励、回合长度、探索率、损失、回放池大小和梯度更新次数等信息，并定期保存模型检查点。基础实验每 50000 步进行一次 20 局开发评估，使用评估均分选取最佳模型。

\subsubsection{低学习率续训与策略保护}
基础模型获得一定控制能力后，项目开展低学习率续训，尝试在已有策略附近继续改进。续训阶段预算为 200000 步，学习率从 $3\times10^{-5}$ 线性下降到 $10^{-5}$，探索率从 0.05 下降到 0.02，每 25000 步评估 50 局。减小学习率有助于缩小单次参数变化幅度，但能否维持已有能力仍取决于新经验分布和价值估计。

现有记录显示，低学习率续训在阶段早期取得最佳开发均分，此后继续训练并未保持这一水平。针对该现象，当前续训配置加入起点评估、策略预热和连续多次开发评估未改进时早停等保护。起点评估先把恢复的模型纳入候选，避免一次较弱的新检查点被当成该阶段的最佳模型；预热期间尚未更新参数，重复评估不会消耗早停耐心；产生新最佳结果时同步保存模型和对应回放池。

上述保护是工程整理后的当前配置。历史低学习率实验完整运行至 200000 步，其结果不能被解释为当前早停配置下重新运行的结果。配置用于说明可复现的训练入口，历史表现则按项目已有记录报告。

\subsubsection{后续恢复实验与评估协议}
在低学习率续训之后，项目又尝试了三步回报恢复和一步扩大回放池两种设置。前者希望加快延迟奖励的利用，后者希望延长历史经验的保留时间。两次实验均未超过恢复起点的开发均分，并提前停止。历史配置快照显示，两项恢复实验采用相同的扩大回放容量；它们从同一个续训最佳模型分别开始，而非彼此顺序衔接。

开发评估在独立环境中进行，使用贪心策略，且不更新参数。训练回调在固定随机性设置下生成连续游戏序列，模型选择所用数据属于开发数据。基础训练每次评估 20 局，续训与恢复实验每次评估 50 局，样本规模不同。

独立复测采用逐局重置的方式，使不同候选模型面对相同的开局集合，共运行 30 局。它与训练回调的连续游戏序列是两套不同协议，各自的均分分别报告，不混合计算。

\reportfigure{images/training_curves.png}{训练与开发评估曲线}{上排为训练回合奖励，下排为贪心开发评估均分；各列独立计步。}{四个实验阶段的实际训练日志。上排为全部回合得分及滚动均值，下排为各评估节点均分；星号标出最佳节点，虚线标出两项恢复实验的早停，浅灰背景表示续训新回放池的预热区间。}{fig:curves}

\subsubsection{探索率与损失日志诊断}
为进一步观察训练信号，图\ref{fig:diagnostics}直接读取两个主要阶段的训练日志，分别绘制探索率和 Huber 损失。日志中的缺失损失值不插值，平滑曲线采用有效日志点的滚动中位数；由于日志记录间隔不完全一致，该窗口不代表固定的决策步数。

从头训练的探索率降至设定下限后，开发均分仍出现明显回落，说明后期回退不能仅用探索率继续下降来解释。低学习率续训的探索率缓慢下降，最佳结果却出现在续训早期，也说明探索安排和策略得分之间没有简单的单调关系。

损失同样不能作为选模的替代指标。基础训练早期的损失较小，后期有所增大，但较小的拟合误差并未对应较强的游戏控制。随着模型访问状态和回放样本变化，TD 目标也会变化，因此损失的尺度和波动必须结合当前经验分布理解。这些日志支持使用独立开发得分选取模型，但尚不足以单独识别回退的具体原因。

\reportfigure{images/training_diagnostics.png}{探索率与损失诊断}{基础训练和低学习率续训的实际探索率、Huber 损失日志。}{两项主要实验的探索率与损失变化。上排为探索概率，下排为对数轴下的原始损失点及滚动中位数；红色虚线标记开发最佳节点。预热期间没有损失记录的区间保留为空白。}{fig:diagnostics}

\subsection{开发结果与行为分析}
各阶段结果及训练曲线来自训练监控与开发评估日志，保留全部实际采样点。为补齐配对分布图，本次固定已有权重，在 CPU 上对两个候选模型分别运行相同开局的 30 局贪心评估；均分复现原记录的 260.27 与 122.40。逐局分数、原始行为帧和观测样本保存在项目的实验数据目录中，这些复测仍属于开发数据。

\subsubsection{各训练阶段的结果}
表\ref{tab:results}给出当前物理版本下 DQN 各阶段的实验结果。表中最佳节点指该阶段开发评估均分最高的节点；后两个恢复实验的最佳模型均为恢复起点，而非新训练得到的模型。

\begin{table}[!htbp]
 \centering\small
 \caption{各阶段开发评估结果}\label{tab:results}
 \begin{tabularx}{\linewidth}{Y P{3.0cm} r r c}
 \toprule
 实验阶段 & 最佳模型所在阶段 & 最佳均分 & 结束均分 & 局数\\
 \midrule
 从头训练 500k & 训练中后期 & 174.00 & 133.30 & 20\\
 一步低学习率续训 200k & 续训早期 & 283.48 & 113.12 & 50\\
 三步回报恢复 & 恢复起点 & 283.48 & 187.78 & 50\\
 一步扩大回放 & 恢复起点 & 283.48 & 216.94 & 50\\
 \bottomrule
 \end{tabularx}
 \vspace{.3em}
 \begin{minipage}{\linewidth}\footnotesize
 注：$1\mathrm{k}=1000$ 个决策步；“局数”为每次评估局数。不同阶段的评估规模不同，表中数值用于呈现各阶段轨迹与模型选择结果。
 \end{minipage}
\end{table}

基础训练在中后期取得 174.00 的最佳均分，结束时下降到 133.30，减少 40.70 分。低学习率续训的最佳均分为 283.48，出现在续训早期；结束时为 113.12，比最佳节点低 170.36 分。两组结果共同说明，增加训练步数并不保证策略持续改善，最终模型与最佳模型需要分别保存。

基础阶段和续训阶段的最高均分来自不同局数的评估，并且续训同时改变了模型起点、训练预算、学习率与探索安排。因此，不能把二者的差值直接归因于降低学习率。更有信息量的比较来自同一协议下的候选模型复测。

\subsubsection{相同开发种子下的候选模型复测}
在相同开局的 30 局开发复测中，当前推荐的最佳模型取得平均分 260.27、中位数 257.00；上一轮最终模型的平均分为 122.40，如表\ref{tab:audit}所示。

\begin{table}[!htbp]
 \centering
 \caption{相同 30 个开发种子上的独立复测}\label{tab:audit}
 \begin{tabularx}{.90\linewidth}{Y r r c}
 \toprule
 候选模型 & 平均分 & 中位数 & 局数\\
 \midrule
 当前推荐最佳模型 & 260.27 & 257.00 & 30\\
 上一轮最终模型 & 122.40 & 101.50 & 30\\
 \bottomrule
 \end{tabularx}
\end{table}

两者平均分相差 137.87 分，相对增加约 112.64\%。由于这次比较使用相同种子集合，它支持在所测开局条件下选择当前推荐模型。最佳模型在 28/30 个种子上得分更高；其得分范围为 36--444 分，仍存在明显的局间波动。箱线图与逐局散点共同呈现这种差异，不能仅凭均值和中位数接近就推断策略稳定。

本次配对复测使用同一低学习率续训阶段保存的最佳检查点与结束检查点。两者均固定权重进行推理，具体模型与逐局分数可由项目中的评估记录核对。

\reportfigure{images/score_distribution.png}{相同开局下的得分分布}{30 局真实复测的配对散点图及箱线图。}{候选模型在相同开发开局上的真实逐局得分。左侧连线对应同一开局，右侧展示箱线图与全部散点，红色菱形为均值。最佳模型在 28/30 局胜出，平均配对差为 137.87 分。}{fig:distribution}

\subsubsection{低分风险与配对差的进一步分析}
平均分提升之外，还需要考察策略是否减少了低分局。表\ref{tab:robustness}根据相同 30 个开发种子的逐局结果补充低分位数与低分局比例。分位数采用线性插值计算；以总分低于 50 分作为描述性阈值，该阈值用于比较两模型的低分尾部，不表示环境中的额外终止条件。

\begin{table}[!htbp]
 \centering\small
 \caption{相同 30 个开发种子上的波动与低分风险}\label{tab:robustness}
 \begin{tabularx}{.92\linewidth}{Y r r}
 \toprule
 指标 & 推荐最佳模型 & 续训最终模型\\
 \midrule
 得分标准差（总体形式） & 103.75 & 90.63\\
 第 10 百分位得分 & 111.90 & 19.90\\
 第 25 百分位得分 & 210.00 & 36.75\\
 低于 50 分的局数 & 1/30 & 9/30\\
 低于 50 分的比例 & 3.33\% & 30.00\%\\
 最低分--最高分 & 36--444 & 3--345\\
 时间截断局数 & 0 & 0\\
 完整通关局数（864 分） & 0 & 0\\
 \bottomrule
 \end{tabularx}
\end{table}

最佳模型的第 25 百分位得分为 210.00，续训最终模型为 36.75；低于 50 分的比例由 30.00\% 降至 3.33\%，相差 26.67 个百分点。这说明在所测开发开局上，提升同时体现在典型得分和低分尾部。最佳模型的绝对标准差仍达到 103.75，且没有一局完整通关，因此上述改进仍伴随局间波动，不能表述为稳定通关。

对每个开局计算配对差 $\Delta_i=R_i^{\mathrm{best}}-R_i^{\mathrm{final}}$，可进一步观察收益是否覆盖全部开局。图\ref{fig:robustness}显示，最佳模型在 28 局中得分较高，但在另外两局中分别少得 198 分和 9 分。均值差为 137.87 分，并不意味着每个开局都改善。对完整回合对反复重采样，得到均值差的百分位 Bootstrap 95\% 区间为 $[95.77,177.43]$ 分。该区间仅描述固定两模型在当前开发样本上的配对差变异，不覆盖独立重复训练的变异，也不替代最终留出测试。

\reportfigure{images/score_robustness.png}{低分风险与配对差}{实际逐局分数的经验累积分布及排序配对差。}{低分风险与逐种子收益。左侧为经验累积分布，虚线标记 50 分阈值；右侧为从小到大排序的最佳模型减最终模型得分，负值对应局部退步，红色虚线为平均配对差。}{fig:robustness}

\subsubsection{训练回退与恢复实验分析}
训练后期的回退首先反映出模型参数更新与策略性能之间并非单调关系。DQN 的学习目标包含目标网络估计，回放池又不断接收新策略产生的数据。即使当前小批量上的损失降低，也可能改变关键状态下的动作排序，进而影响之后整条轨迹。一次错误移动造成失球，还会减少后续继续获得奖励的机会。

回放分布变化是可能的解释之一。当回放容量有限时，较早阶段的有效经验会逐步被新经验替换；若新策略访问的状态变窄，网络可能减少对某些关键接球状态的训练。另一方面，自举估计误差、探索变化和不同开局下的轨迹差异也可能影响结果。现有汇总统计不能单独识别这些因素的贡献，因而回退原因仍属于需要通过日志和控制实验检验的解释。

三步回报恢复结束均分为 187.78，一步扩大回放结束均分为 216.94，均低于起点的 283.48。这说明在已执行设置下，两种恢复方案都没有形成超过起点的新最佳策略。三步目标引入更多真实奖励，并不保证在当前数据分布上优于一步目标；扩大回放也会同时改变经验新旧比例，不能简单等同于增强稳定性。

早停和最佳模型保存的作用是保留已有成果并减少继续退化阶段的计算支出。它们改变的是训练管理和模型选择方式，并没有使每一步更新都具备性能保证。当前选择依据应是独立评估结果，而不是模型文件的生成时间或累计训练步数。

\subsubsection{实时运行与行为分析}
实时游戏画面适合补充统计分数背后的行为解释。对同一局游戏，应按时间顺序记录小球下降、挡板移动、接触瞬间和反弹后的连续画面，同时保留模型、种子、分数及生命信息。仅有单帧画面难以判断运动方向，也不能说明一次接球是主动预测还是偶然重合。

\reportfigure{images/behavior_sequence.png}{模型行为连续帧}{同一模型、同一局中的成功接球与失球连续帧。}{同一局中的接球与失球对照，每排按相同时间间隔顺序采样。画面裁取下方接球区域，黄圈标记小球，蓝箭头标记挡板位移；子图保留实际得分与剩余生命，用于辅助判断运动变化。}{fig:behavior}

行为分析重点关注挡板与小球的相对运动。图\ref{fig:behavior}取自推荐最佳模型在同一局游戏中的两个片段。成功片段中，挡板在小球下降期间向接触位置移动，小球随后向上反弹。失球片段中，小球在右侧接近挡板高度时横向位置超出挡板覆盖范围，随后落到挡板下方，剩余生命减少。后续向左移动未能挽回这次失球。这些连续帧说明该模型仍会在局部轨迹上失误；单个成功片段不能证明其已学会精确落点预测，也不能据此推断所有失球的原因。

DQN 开发流程最终保留低学习率续训的最佳检查点。其与 PPO 的冻结模型测试统一见第\ref{sec:comparison}节，开发复测均分不与最终测试均分混合计算。

\FloatBarrier

\clearpage
\section{PPO 方法与实验}
\label{sec:ppo}
本节按与 DQN 相同的层次介绍 PPO，重点解释同策略更新、实现变体、开发选模与配对对照。除已有 GAE 消融外，补充 KL 学习率反馈的三次配对续训、历轮改动的收益区间及控制过程诊断。历轮开局集合不同，分数分别按各轮协议解释；DQN 与 GAE 阶段冻结候选的同开局 50 局结果集中见第\ref{sec:comparison}节。

\subsection{算法原理}
\subsubsection{PPO 的策略与价值学习}
与 DQN 估计各动作的 Q 值不同，近端策略优化（PPO）直接学习动作分布 $\pi_\theta(a\mid s)$，并用状态价值网络 $V_\phi(s)$ 辅助评价动作。策略网络输出四个动作的 logits，经 softmax 得到概率；训练时按概率采样，确定性评估时选择概率最大的动作。价值网络输出一个标量，估计当前状态下的未来折扣回报。

PPO 属于同策略方法。它先用旧策略收集一段 rollout，记录状态、动作、奖励、价值预测与旧动作对数概率，再在这批数据上进行若干轮小批量更新。该批更新后继续采集新轨迹，不把任意旧回放直接代入 PPO 概率比目标。它不使用 DQN 的目标 Q 网络或 $\varepsilon$-greedy 探索；探索由动作分布和熵奖励共同调节。

\subsubsection{GAE 优势估计与信用分配}
定义价值时间差分误差与广义优势估计为
\begin{align}
 \delta_t^V &= r_t+\gamma(1-d_t)V_{\phi_{\mathrm{old}}}(s_{t+1})-V_{\phi_{\mathrm{old}}}(s_t),\\
 \widehat A_t &= \sum_{l=0}^{K-1}(\gamma\lambda)^l\delta_{t+l}^V,
 \qquad \widehat R_t=\widehat A_t+V_{\phi_{\mathrm{old}}}(s_t).
\end{align}
其中 $d_t$ 只在真正终止时为 1；$K$ 为从 $t$ 到本局终止、时间截断或当前 rollout 边界前的有效转移数，不跨越环境重置累加优势。真正终止不自举，时间截断则用实际终止观测的价值自举。$\lambda$ 调节偏差和方差，$\widehat R_t$ 是固定的价值回归目标。训练小批量内的优势标准化不改变评估的原始得分。

GAE 对照阶段固定 $\gamma=0.995$，仅比较 $\lambda=0.98$ 与 0.99。权重衰减尺度 $1/(1-\gamma\lambda)$ 约从 40 增至 67 个决策步，有望将较远的击砖奖励联系到更早动作，但也可能增加优势方差。这是权重尺度，不是超过该步数就完全忽略回报；价值自举仍参与长期估计。

\subsubsection{概率比裁剪与训练目标}
新旧策略对已采样动作的概率比为
\begin{equation}
 \rho_t(\theta)=\frac{\pi_\theta(a_t\mid s_t)}{\pi_{\theta_{\mathrm{old}}}(a_t\mid s_t)}.
\end{equation}
PPO 的裁剪代理目标及本项目最小化的总损失为
\begin{align}
 L^{\mathrm{clip}}(\theta)&=\mathbb E_t\!\left[\min\!\left(\rho_t\widehat A_t,
 \operatorname{clip}(\rho_t,1-\epsilon,1+\epsilon)\widehat A_t\right)\right],\\
 \mathcal L(\theta,\phi)&=-L^{\mathrm{clip}}(\theta)
 +c_V\mathbb E_t\!\left[\ell_V\bigl(\widehat R_t-V_\phi(s_t)\bigr)\right]
 -c_H\mathbb E_t\!\left[\mathcal H(\pi_\theta(\cdot\mid s_t))\right].
\end{align}
其中 $\epsilon$ 为概率比裁剪范围，$c_V$ 和 $c_H$ 控制价值回归和熵奖励。裁剪限制继续放大某些概率变化的收益，但不保证每次更新都提高游戏得分；项目另设近似 KL 保护。独立梯度裁剪和稳健价值回归属于 PPO 的实现变体。

\subsection{模型实现与训练管理}
\subsubsection{PPO 策略与价值双网络}
初始标准 PPO 共享卷积特征，后续改进版本将策略与价值特征提取器分开。每个分支均采用 NatureCNN：三层卷积依次使用 32 个 $8\times8$、64 个 $4\times4$ 和 64 个 $3\times3$ 卷积核，步长为 4、2、1，展平后连接 512 维特征，与表\ref{tab:network}的特征提取部分一致；策略头为 $512\rightarrow4$ logits，价值头为 $512\rightarrow1$ 标量。拆分时复制原权重，使迁移前后的动作概率和价值预测相同；参数组改变后重置 Adam，匹配对照做相同重置。推理只使用策略分支。

Balanced PPO 保留标准 rollout、GAE、概率比裁剪和优势标准化，对策略与价值参数分别按范数 0.5 和 1.0 裁剪。价值损失为 $2\times\mathrm{Huber}(\delta=10)$：小误差区与 MSE 相同，大误差区减弱极端残差影响。网络分开不自动意味着梯度预算独立；裁剪缩放也不能直接解释为 Adam 参数步长的等比例变化。

具体地，令价值残差 $e=\widehat R-V_\phi(s)$，后续 GAE 与 KL 对照均采用
\begin{equation}
 \ell_V(e)=\begin{cases}
 e^2, & |e|\leq 10,\\
 20|e|-100, & |e|>10.
 \end{cases}
\end{equation}
因此单样本损失对残差的导数绝对值最多为 20；这与参数梯度范数裁剪作用于不同层面。GAE 与 KL 对照未启用价值预测裁剪，避免将其与策略概率比裁剪混淆。

\reportfigure{images/ppo_actor_critic.png}{PPO 策略与价值网络}{相同图像输入分别送入独立 CNN，输出动作分布与状态价值。}{PPO 改进版本的策略/价值结构。训练时采样动作并拟合 GAE 目标，确定性评估只选择策略概率最大的动作。}{fig:ppo-network}

\subsubsection{训练循环与模型更新}
PPO 的更新循环按完整 rollout 组织，后续对照的共同流程如下：
\begin{enumerate}
 \item 用当前策略在 16 个环境中各采集 256 个决策步，保存旧动作概率和旧价值预测，共得到 4096 条转移。
 \item 在局末或 rollout 边界正确自举，计算 GAE 和固定 returns；后续 epoch 不重算旧策略概率或优势目标。
 \item 打乱这批数据，以 256 条转移为一个小批量，最多优化 4 个 epoch；小批量内标准化优势，分别计算策略、价值和熵项。
 \item 分别裁剪策略与价值梯度，按退火后的学习率更新 Adam；若小批量近似 KL 超过目标的 1.5 倍，则提前结束本轮剩余更新。
 \item 在整批原 rollout 上记录更新后的 KL 和裁剪比例，随后采集新轨迹；按开发协议评估并保留候选模型。
\end{enumerate}
一批轨迹样本内的重复优化属于同一轮 PPO 更新，不是跨轮长期回放。训练预算约为 50 万个决策步，按完整 rollout 批次完成采样，因此实际步数会略高于计划预算。

\subsubsection{续训与变体状态保存}
PPO 续训恢复策略、价值与优化器状态，重新采集 rollout，不恢复历史回放。阶段衰减按新增步数计时，各变体均保存相应元数据。DualRate 的额外价值拟合保持原优势、returns 和旧动作概率，不更新策略权重及其 Adam 状态。KL 反馈变体另外保存平滑 KL、学习率倍率和控制计数，只有显式开始新的优化器阶段时才重置控制状态。

PPO 训练入口为 \texttt{train\_ppo.py}，独立裁剪、双学习率与 KL 反馈封装为更新变体。逐局开发种子与候选选择使用共享模块，以保证不同变体的评估口径一致。

\subsection{训练设置与开发选模}
\subsubsection{PPO 训练阶段与单因素对照}
标准 PPO 首轮从随机初始化训练约 50 万步；之后依次研究独立策略/价值 CNN 与退火、大 rollout、长回报、辅助奖励、独立梯度裁剪和稳健价值损失。辅助奖励曾设置接球奖励与失球惩罚，改变训练目标，未作为最终比较方案。未获得收益的实验同样保留。GAE 同起点对照之后，又开展 KL 反馈的多次配对续训；不同测试开局上的历轮分数不拼成提升曲线。

GAE 阶段用新验证集选出的同一个 Balanced 检查点作为起点，采用相同的随机性设置，两组均重置 Adam。主要参数见表\ref{tab:ppo-hyper}，仅 GAE $\lambda$ 不同，均按相同预算完成训练。多个并行环境由同进程向量环境采样，不代表多次独立训练。

\reportfigure{images/ppo_training_curves.png}{PPO 训练和开发诊断}{同一模型起点的两组 GAE 设置，展示真实开发得分与更新日志。}{PPO GAE 对照的实际日志。开发均分与中位数使用相同 32 个逐局种子；训练均分含随机动作，整轮 KL 在固定 rollout 上更新后测量，不能与开发回报混为一谈。}{fig:ppo-curves}

\begin{table}[!htbp]
 \centering\small
 \caption{PPO 的 GAE 阶段共同训练设置}\label{tab:ppo-hyper}
 \begin{tabularx}{.95\linewidth}{Y Y}
 \toprule
 参数 & 设置\\
 \midrule
 单组新增预算 & 约 50 万个决策步\\
 环境数 $\times$ 每环境 rollout & $16\times256=4096$\\
 批大小 / 优化 epoch & 256 / 4\\
 学习率 & $2\times10^{-5}\rightarrow5\times10^{-6}$\\
 折扣因子 / GAE lambda & 0.995 / 0.98 或 0.99\\
 clip range / KL 目标 & $0.06\rightarrow0.03$ / 0.0075\\
 熵系数 / 价值损失权重 & $0.0003\rightarrow0.0001$ / 0.25\\
 策略 / 价值梯度范数上限 & 0.5 / 1.0\\
 价值损失 / Adam & $2\times\mathrm{Huber}(\delta=10)$ / 两组均重置\\
 随机性设置 / 开发评估 & 两组保持一致 / 定期评估 32 局\\
 \bottomrule
 \end{tabularx}
\end{table}

\Needspace{8\baselineskip}
\subsubsection{开发数据与候选选择}
起点选择、训练开发评估与最终测试使用相互分离的开局集合；开发环境按局重置随机性设置，不依赖前局的随机数消耗；最终测试对各候选统一运行 50 局。最终优先选开发均分最佳，只有中位数最佳模型的均分达到其 95\% 且中位数更高时才替换。该阈值是预先约定的取舍，不是显著性检验。GAE 阶段两组均执行完整预算，稳定候选均未满足最终均分约束，实际选中均分最佳模型。

\subsubsection{KL 反馈与三次配对续训设置}
KL 阶段从 GAE 开发流程选中的 $\lambda=0.98$ 模型冻结副本开始，保留原始奖励和上述独立网络、Huber 损失及梯度裁剪。采用三种不同的训练随机性设置，每种设置各运行反馈组和固定计划对照，组成三次配对续训。每次新增约 50 万个决策步，实际均为 503808 步；两组均重置 Adam，执行完整预算。每个检查点评估相同的 32 局开发开局，最终所选模型各做 50 局独立测试。此阶段另用新的开发和测试开局集合，不与第\ref{sec:comparison}节的原统一测试混合。

基础学习率仍由 $2\times10^{-5}$ 退火到 $5\times10^{-6}$，仅反馈组的策略分支乘以可调倍率：
\begin{equation}
 \eta_{\pi,k}=m_k\eta_{\mathrm{base},k},\qquad
 \eta_{V,k}=\eta_{\mathrm{base},k},\qquad
 m_k\in[0.5,2].
\end{equation}
控制器观测整轮更新后的采样 KL，以 0.2 为新观测权重进行指数平滑。反馈目标为 $5\times10^{-4}$，平滑值低于目标的 $1/1.5$ 时将下一轮倍率乘以 1.1，高于目标的 1.5 倍时除以 1.1；瞬时 KL 超过上界时优先下调，最后限制倍率范围。对照使用同一实现关闭反馈，倍率恒为 1，价值学习率及其余设置一致。该反馈目标与小批量 KL 保护参数 0.0075 作用不同：前者调整下一轮策略学习率，后者限制当前更新。

\subsection{优化结果与负结果分析}
标准 PPO 在首轮 20 局测试中均分为 35.00；稳定续训技巧在新 30 局上为 43.33，未超过单纯续训对照的 53.00。长回报原奖励模型在另一组 50 局上为 88.54，辅助奖励为 57.34；Balanced 在后续 50 局上为 100.76、标准更新对照为 81.84，但配对均分差的区间跨零。各轮测试开局不同，只记录实验过程，不直接计算历轮提升率。

双学习率实验提高价值学习率并增加额外拟合。在已记录的更新中，96.9\% 的额外拟合降低训练 MSE，相对降幅中位数约 19\%，但按开发均分选中的仍为起点。独立测试均分 128.12 与起点相同，Balanced 对照为 134.00。更低训练误差并未自动转化为更好策略。

该阶段稳定候选的开发均分为 158.53，起点为 158.59，两者十分接近；中位数却由 80.5 升至 132.5。这促使项目在新验证集上复核候选并增加均分约束的稳定性选择，而不是按最终测试重新挑模型。随后由开发流程选中的 Balanced 检查点成为 GAE 两组的共同起点。

GAE 阶段主候选由开发均分选定为 $\lambda=0.98$；$\lambda=0.99$ 作为补充消融保留。统一测试同时纳入两组的共同起点，具体分数、配对区间和截断情况在第\ref{sec:comparison}节集中分析。

\subsubsection{多次配对续训与选模稳定性}
图\ref{fig:ppo-multiseed}逐次展示 KL 对照的开发均分和中位数，保留全部实际评估节点，不把三次训练合成一条平滑的提升曲线。各次重复共享模型起点、游戏规则和开发开局，但训练采样的随机性不同；同一次重复的反馈组和对照使用相同训练随机性设置。

\reportfigure{images/ppo_multiseed_curves.png}{PPO 多次配对续训}{三次重复的开发均分、中位数和最终选中检查点。}{KL 反馈与固定计划的三次配对续训。上排为每次 32 局开发均分，下排为中位数；星号表示按开发规则最终选中的检查点，空心菱形表示均分最佳节点，二者可能重合。每列独立展示一次训练重复，未用跨重复置信带代替个体曲线。}{fig:ppo-multiseed}

曲线显示不同重复具有不同的优势阶段。重复一的反馈组在中途回落后恢复，重复二的反馈组开发均分在后期较高，重复三则在较早节点达到自身最佳，而固定计划对照在后期继续改善。由此不能把某一次、某个节点的领先解释为稳定的算法优势，也不能假定更多训练必然提高得分。

选模还需要同时考虑均分和典型表现。重复一反馈组的开发均分最高为 234.53，但该节点中位数为 185.50；依据预先约定的约束，最终选中开发均分 231.56、中位数 203.00 的稳定候选。其余重复同样依据开发数据选择，不因独立测试结果更换检查点。三次重复是同一冻结起点下的续训实验，尚不能代表从随机初始化开始的三次独立训练。

\subsubsection{改进措施的配对消融}
图\ref{fig:ppo-ablations}用“改动减同轮对照”的配对均分差汇总历轮观察。每一行的两组策略面对相同开局，零线表示样本均分相同；区间用于展示差异的不确定性。由于各行的模型起点、测试集合与训练历史不同，图中数值不构成跨轮算法排名。稳定续训和独立裁剪/Huber 均含组合改动，不能据此单独归因某一个部件；双学习率组还因早停具有不同实际预算，且最终保留起点模型，尤其不能视为纯粹的价值拟合单因素收益。

\reportfigure{images/ppo_ablation_effects.png}{PPO 改进措施消融}{各轮相对自身对照的配对收益及不确定性。}{各轮改动相对同轮对照的配对均分差。前五行展示固定策略的回合配对区间，最后一行为 KL 三次续训的训练与游戏两维配对区间；灰点表示三次 KL 续训各自的均分差。不同区间的统计对象不同，不能把跨零与否用作统一的算法优劣检验。}{fig:ppo-ablations}

在各轮固定策略样本中，辅助奖励相对原始奖励的均分差为 \ShapingAblationDiff，区间为 $[\ShapingAblationLow,\ShapingAblationHigh]$，该轮未支持其更优。稳定续训组合、独立裁剪/Huber、双学习率组合和 GAE 调整的区间均跨零，当前证据不足以认定它们具有稳定的额外收益。保留这些结果比只展示有利的中间节点更能反映实际优化过程。

KL 阶段三次续训的测试均分分别取平均，反馈组为 \KLAdaptiveMean，对照组为 \KLControlMean；逐次配对差分别为 $+\KLRepeatOneDiff$、$\KLRepeatTwoDiff$ 和 $+\KLRepeatThreeDiff$。总体配对差为 \KLAblationDiff，按训练重复对和共同游戏开局两个维度进行配对重采样，得到 95\% 区间 $[\KLAblationLow,\KLAblationHigh]$。三次重复中既有提高也有退步，区间跨零，尚不能确认 KL 反馈具有稳定优势。

这里每种设置虽然共测试 150 局，但只有三次续训重复，不能把 150 局当成 150 次独立训练。前几轮区间仅描述固定策略在回合样本上的条件差异；KL 区间额外考虑当前三个续训重复，仍受重复数量少及共享起点的限制。原始历史区间按已有统计保留，新补充的辅助奖励配对区间与 KL 两维区间均使用 20000 次重采样。

\subsubsection{KL 控制过程与更新诊断}
图\ref{fig:ppo-kl-diagnostics}将倍率、实际策略学习率、整轮 KL、平滑 KL、概率比越界比例及裁剪范围按日志记录位置对齐。三次重复全部展示，颜色区分两种设置，线型区分重复。优化指标的日志输出位置与开发评估节点并非同一采样时点，因此这些曲线用于诊断更新过程，不作逐点因果归因。

\reportfigure{images/ppo_kl_diagnostics.png}{KL 控制诊断}{学习率、完整更新后 KL、反馈阈值与概率比裁剪调度。}{KL 反馈的实际运行日志。平滑 KL 子图的阴影表示反馈阈值带；倍率和平滑值子图展示反馈组，倍率为 1 的水平线表示固定计划对照。概率比越界比例定义为 $|\rho-1|>\epsilon$ 的样本比例，不等同于实际被裁剪的梯度贡献。}{fig:ppo-kl-diagnostics}

\FloatBarrier
控制器确实改变了更新幅度：三次反馈组在前期接近固定计划，随后逐步提高倍率，并在后半段达到 2 倍上限；因此阶段末实际策略学习率约为 $10^{-5}$，对照约为 $5\times10^{-6}$，价值分支仍沿基础计划退火。最后约 10 万步中，反馈组整轮采样 KL 的中位数约为 $1.6\times10^{-4}$，对照约为 $8.5\times10^{-5}$；概率比越界比例分别约为 6\% 和 4\%。这些观察支持反馈已生效，但不能把更大的 KL 或越界比例直接解释为更好的控制。

还可以观察到，基础学习率和裁剪范围持续下降，而反馈目标保持固定。后期平滑 KL 低于阈值带时，倍率已经达到上限，控制器无法继续补偿退火。这提示下一轮可分别研究反馈目标的阶段调度或更保守的倍率上限；当前曲线尚不足以认定调度差异就是个别重复退步的原因。结合开发曲线与独立测试的波动，应继续将参数变化、候选选择和最终得分分开评价。

\FloatBarrier

\clearpage
\section{统一测试与方法比较}
\label{sec:comparison}
\subsection{冻结候选与统一测试协议}
本部分固定已有 DQN 推荐模型、GAE 对照的两个 PPO 模型及其共同起点，均在相同物理规则、原始奖励、完整五生命游戏和 $4\times84\times84$ 输入下，用 CPU 确定性动作运行相同开局的 50 局测试。随机基线在相同任务和开局设置下均匀采样动作。所有结果均来自逐局实际运行，不由汇总分数反推。

主比较采用开发均分更高的 PPO $\lambda=0.98$（196.00，高于 0.99 的 171.09）；0.99 作为补充消融同时报告，而不根据其测试均分较高事后更换主候选。DQN 推荐模型仍为原开发流程选出的低学习率续训最佳检查点。

\subsection{统一测试结果与配对比较}
\begin{table}[!htbp]
 \centering\small
 \caption{相同 50 个测试种子下的固定模型与随机基线}\label{tab:final-comparison}
 \begin{tabularx}{\linewidth}{Y r r r c c}
 \toprule
 固定模型 & 均分 & 标准差 & 中位数 & 低于50分 & 截断\\
 \midrule
 随机策略 & \RandomMean & \RandomStd & \RandomMedian & \RandomLow/50 & \RandomTimeout/50\\
 DQN & \DQNMean & \DQNStd & \DQNMedian & \DQNLow/50 & \DQNTimeout/50\\
 PPO 共同起点 & \SourceMean & \SourceStd & \SourceMedian & \SourceLow/50 & \SourceTimeout/50\\
 PPO $\lambda=0.98$ & \PPOControlMean & \PPOControlStd & \PPOControlMedian & \PPOControlLow/50 & \PPOControlTimeout/50\\
 PPO $\lambda=0.99$ & \PPOGaeMean & \PPOGaeStd & \PPOGaeMedian & \PPOGaeLow/50 & \PPOGaeTimeout/50\\
 \bottomrule
 \end{tabularx}
\end{table}

对相同种子计算 $\Delta_i=R_i^{\mathrm{DQN}}-R_i^{\mathrm{PPO},\lambda=0.98}$。配对均分差为 \DQNPPODiff\ 分，DQN 胜/平/负为 \DQNPPOWins/\DQNPPOTies/\DQNPPOLosses。

\reportfigure{images/dqn_ppo_comparison.png}{DQN 与 PPO 的统一测试}{50 个固定种子的得分分布及 DQN 减 PPO 主候选的配对差。}{统一测试协议下的实际结果。左侧展示全部逐局分数与箱线图，菱形表示均值；右侧为按大小排序的 DQN 减 PPO $\lambda=0.98$ 得分差，负值表示相应开局 PPO 更高。}{fig:method-comparison}

对完整回合对反复重采样，得到配对均分差的 95\% bootstrap 区间 $\DQNPPOCI$。这一比较衡量固定模型在当前测试开局上的条件差异，不覆盖独立训练波动。截断表示达到统一时间上限，不等于通关；局长增加也不必然表示单位交互得分效率提高。

\subsection{PPO 的 GAE 消融}
表\ref{tab:final-comparison}中的共同起点与两组 PPO 使用相同 50 个种子，提供同起点后续训练的评价。共同起点的均分为 \SourceMean、中位数为 \SourceMedian，第 25 百分位为 \SourcePxxv；$\lambda=0.98$ 和 0.99 的第 25 百分位分别为 \PPOControlPxxv\ 和 \PPOGaePxxv。两组新增步数与训练设置相同，只有 GAE $\lambda$ 不同，主候选仍按开发数据选择。

0.99 相对 0.98 的配对均分差为 \GaeControlDiff，95\% bootstrap 区间为 $\GaeControlCI$。应同时观察区间、低分尾部和截断，不能只因一个样本均分更高就认定该设置普遍更优；两组各只有一个训练种子。

\subsection{学习机制与工程取舍}
两种方法处理相同图像与动作空间，但优化对象和数据利用方式不同。DQN 通过 Q 值排序间接产生动作，PPO 直接更新策略分布；两者的差异不能只归结为某一个超参数。

\begin{table}[!htbp]
 \centering\small
 \caption{本项目 DQN 与 PPO 的机制比较}\label{tab:method-comparison}
 \begin{tabularx}{\linewidth}{P{2.3cm}YY}
 \toprule
 方面 & DQN & PPO\\
 \midrule
 优化对象 & 四动作 Q 值的 TD 回归 & 裁剪策略目标与状态价值回归\\
 数据利用 & 离策略；长期回放、随机抽样复用 & 同策略；当前 rollout 内多轮更新\\
 稳定机制 & 目标 Q 网络、回放与 Huber & 概率比裁剪、优势标准化、KL 保护\\
 探索方式 & $\varepsilon$-greedy & 概率采样与熵奖励\\
 最终网络 & 在线 Q 网络，另有目标网络 & 独立策略/价值 CNN，推理只用策略\\
 续训状态 & 模型与回放分开保存 & 模型与优化器保存，重新采集 rollout\\
 主要取舍 & 回放复用经验，图像回放内存较大 & 轨迹更新及时，跨 rollout 样本复用受限\\
 \bottomrule
 \end{tabularx}
\end{table}

PPO 的裁剪没有消除训练波动，DQN 的经验回放也没有保证继续训练必然改善。两个方法都需要起点评估、检查点保留及独立开发集选模；训练损失不宜跨算法直接比较。

\subsection{比较结论与适用范围}
DQN 的第 25 百分位为 \DQNPxxv，高于主 PPO 的 \PPOControlPxxv；低于 50 分的局数分别为 \DQNLow/50 与 \PPOControlLow/50。结合均分、中位数与配对区间，这组固定模型的差异也体现在典型表现和低分尾部。

PPO 两组分别有 \PPOControlTimeout/50、\PPOGaeTimeout/50 局达到时间上限，DQN 为 \DQNTimeout/50；所有候选均未取得 864 分完整通关。较长轨迹不能被解释为更强的通关能力，训练预算和重复性带来的结论边界见第\ref{sec:conclusion}节。

\FloatBarrier

\clearpage
\section{总结与局限}
\label{sec:conclusion}
\subsection{项目工作总结}
本项目在统一 Breakout 环境中实现了 DQN 与 PPO 的图像控制、训练、开发选模和实时展示。两者共用动作、奖励、回合终止及观测接口；DQN 以回放和目标网络学习动作价值，PPO 以 GAE 与裁剪目标优化动作分布。日志、检查点和逐局分数使实验可以复核。

DQN 低学习率续训得到较好的中间模型，继续训练和后续恢复却未进一步提高得分。PPO 的额外价值拟合降低了训练误差，但未改善开发所选策略；提高 GAE $\lambda$ 的额外收益也未得到区间支持。新增 KL 对照虽实际提高了策略学习率和整轮 KL，三次续训的测试收益仍有正有负，总体区间跨零。两条链路的共同经验是保留负结果、依据开发集选模，并用独立测试检查策略表现。

\subsection{局限与后续验证}
原 DQN 的配对复测仍属于开发数据。本次在冻结模型后，使用留出的共同开局完成同协议的 50 局测试及随机基线，已在第\ref{sec:comparison}节集中比较；未依据这些测试更新权重或更换主比较候选。结果仅支持固定模型在当前规则下的比较，不能证明稳定通关或跨环境泛化。

本实验并非同预算、同初始化的算法排名：DQN 与 PPO 使用不同的多阶段训练历史和超参数搜索，既不能用单次测试差推断普遍优劣，也不能从当前耗时推断公平的样本效率。更强结论需要统一从头训练预算、多个独立训练种子和未用于调参的新测试集。当前最可靠的结论是两个已训练系统在同一任务上的可复核表现，以及训练、选模和稳定性取舍的实测差异。

后续需对入围设置进行多个独立训练种子的重复实验，并保留未用于调参的测试集。除平均分外，还应报告标准差、中位数、低分位数、低分比例、通关率和时间截断比例。对同一模型增加开局种子，不能替代独立重复训练。

对于学习率、回报跨度和回放容量的作用，可以从同一模型起点出发，固定其余设置并分别改变一个因素。这样能够进一步区分模型起点、额外交互预算和参数设置对结果的影响。当前恢复实验的负结果应保留为已有观察，而不是在缺少新数据时替换为预期提升。

\FloatBarrier

\end{document}
